\section{Introduction}

\IEEEPARstart{T}{he} human visual system is capable of solving under-constrained inverse rendering problems. Many different combinations of shape, materials and lighting can produce the same observed image \cite{belhumeur_bas-relief_1999}. Hence, in order to estimate the shape, material properties and colour of surfaces and objects in a scene, the human visual system draws on priors that seek the most likely explanation. These include strong priors over the space of possible illuminations \cite{murray2019visual}. For example, there is wide evidence that humans exploit a lighting-from-above or increasing-luminance-with-elevation prior \cite{hill_independent_1993, thomas_interactions_2010}. It has also been shown that humans perform inverse rendering tasks better under complex natural illumination while performance degrades when the illumination statistics are not representative of the real world \cite{fleming2003real}. Interestingly, it seems that these illumination priors are, at least partially, learnt \cite{granrud1985infants,thomas2010interactions} and can be updated from experience \cite{adams2004experience}.

Compact but expressive lighting representations play an essential role in graphics, enabling realistic lighting effects at real-time frame rates \cite{ramamoorthi_efficient_2001, tsai_all-frequency_2006, wang_all-frequency_2009, ng_all-frequency_2003, green_spherical_2003} and in computer vision, enabling scene relighting \cite{yu_outdoor_2021, yu_self-supervised_2020}, face relighting \cite{wang_face_2009, shu_neural_2017, egger_occlusion-aware_2018, sengupta_sfsnet_2018} and object insertion \cite{wang_learning_2021, song_neural_2019, li_inverse_2020}. Real-world illumination is highly complex and variable, with a very high dynamic range, and is therefore inherently challenging to represent. However, real-world illumination does contain statistical regularities \cite{dror_statistical_2004}, particularly for outdoor, naturally lit scenes. The lighting-from-above prior holds in a world where the strongest illumination source is the sun or skylight, which also produce only a limited range of colours. In addition, illumination environments have a canonical up direction (vertical axis aligns with gravity) but arbitrary horizontal rotation (any rotation about the vertical is equally likely). We also desire vision systems to have exposure invariance (i.e.~a human would infer the same inverse rendering result regardless of pupil dilation) and so an illumination environment could be encountered with any absolute scale. These regularities and geometric symmetries can significantly restrict the space of possible illuminations to constrain inverse problems or enable the synthesis of realistic lighting.

Given this, it is surprising that statistical illumination priors have been almost completely ignored in computer vision, with the vast majority of inverse rendering techniques allowing arbitrary illumination within their chosen representation space. In this paper we attempt to replicate the sort of illumination prior used by the human visual system. This entails learning a statistical characterisation of the illumination environments likely to be encountered in the real world but also incorporating invariances and symmetries by construction. Specifically, we desire a representation of natural illumination environments that exhibits the following features:
\begin{itemize}
    \item \emph{Generative:} A generative model that captures the statistical regularities of natural illumination with a well-behaved latent space within which we can optimise to solve inverse problems.
    \item \emph{Compact:} Reduces the dimensionality of inverse problems while preserving high-frequency lighting effects that are important for non-Lambertian appearance.
    \item \emph{Rotation-Equivariant:} Respect the canonical orientation, i.e.~any rotation of an environment about the vertical should be equally likely and equally well represented.
    \item \emph{Scale-Invariant:} Any scaling of the exposure of an environment should be equally likely and equally well represented.
    \item \emph{Statistical Prior:} Provides a prior to regularise inverse problems, or that can be sampled from for synthesis, only generating plausible illumination environments.
    \item \emph{HDR:} Correctly handle HDR quantities essential for realistic rendering and the reproduction of natural light.
\end{itemize}

\subsection{Contributions}

We introduce RENI++ - A Rotation-Equivariant Natural Illumination model. In so doing, we make the following key contributions:
\begin{itemize}
\item[--] An extension of Vector Neurons to a rotation-equivariant neural field representation for spherical images, optionally restricted to rotations about the vertical axis.
\item[--] A variational autodecoder architecture for a generative model of spherical signals.
\item[--] The first natural, outdoor HDR illumination model.
\item[--] Evaluated in an inverse rendering task showing significant performance improvements over other lighting representations.
\end{itemize}
We choose to model scene radiance, i.e.~environment lighting, directly as opposed to pre-integrated lighting with a particular BRDF. This makes our model more general since it can be used with arbitrary BRDFs at inference time or even for tasks other than rendering, such as to constrain shape from specular flow.

This work is an extension of our earlier RENI model \cite{gardner2022rotation}. Compared to our earlier work, RENI++ adds: 
\begin{itemize}
\item[--] A scale-free loss enabling a large increase in generalisation capabilities on previously out-of-distribution HDR environments.
\item[--] Transformer-decoder architecture with positional encoding, replacing the previous SIREN network, able to capture higher frequency details.
\item[--] An invariant representation that scales as \(O(n)\) rather than \(O(n^{2})\) with the size of latent space by replacing the Gram-Matrix with the VN-Invariant layer \cite{deng_vector_2021}.
\item[--] We generalise the $SO(2)$ formulation to an arbitrary rotation axis in Section \ref{sec:SO2}.
\item[--] A new NeRFStudio-based \cite{nerfstudio} implementation with an order-of-magnitude speedup in training time. 
\item[--] A new data normalisation and sampling procedure resulting in smoother training and faster convergence.
\end{itemize}
\subsection{Overview}

In Section \ref{sec:related_work} we review prior work in the areas of lighting representations, illumination priors, neural fields and invariances. In Section \ref{sec:RECSNF} we introduce our rotation-equivariant, conditional neural field representation for spherical signals. In Section \ref{sec:RENI} we describe how this framework is applied to the task of learning a natural illumination prior, including the use of a scale-invariant loss. In Section \ref{sec:eval} we evaluate our model for generalisation, interpolation, completion, inverse rendering and LDR to HDR. We also perform an ablation study of our design decisions.
